\title{Matemáticas 1º ESO}

\href{/teaching/sec-teaching/}{\icon{fas fa-rotate-left} Back to Secondary Education.}

\maketitle

\section*{(B1) Números y operaciones}

\subsection*{Conteo}

\begin{itemize}
  \item Estrategias de recuento sistemático.
\end{itemize}

\subsection*{Cantidad}

\begin{itemize}
  \item Estimaciones con precisión adecuada.
  \item Números enteros para expresar cantidad.
  \item Números fraccionarios y decimales para expresar cantidad. Diferentes formas de representación (recta numérica)
  \begin{itemize}
    \item Gairín y Sancho. \emph{Números y algoritmos}, Capítulos 8--10.
  \end{itemize}
  \item Valor absoluto de un entero.
  \item Clasificación de los números reales. \textcolor{red}{\textbf{(????)}}
\end{itemize}

\subsection*{Operaciones}

\begin{itemize}
  \item Estrategias de cálculo mental con naturales.
  \item Jerarquía de las operaciones. Interpretación de las operaciones.
  \begin{itemize}
    \item Jugar con los paréntesis y replicar distintas situaciones de la vida real. El orden refuerza el sentido de operador de suma o resta.
    \item El problema de los cuatro cuatros.
    \item Descomposición en suma de un número que maximiza el producto de dicha descomposición.
    \item \href{/problema_combinadas_1eso.jpg}{Problema de combinadas}
  \end{itemize}
  \item Relaciones inversas entre operaciones.
  \begin{itemize}
    \item Conexión con el punto anterior.
  \end{itemize}
  \item Potencias de exponente entero.
  \item Números fraccionarios y decimales
  \begin{itemize}
    \item Necesario trabajar el número racional desde la medida (B2)
    \item Gairín y Sancho. \emph{Números y algoritmos}, Capítulos 8--10.
    \item \href{https://tierradenumeros.com/post/hilo-multiplicacion-de-fracciones-unpacked/}{Multiplicación}
    \item \href{https://tierradenumeros.com/post/hilo-division-fracciones-unpacked/}{División}
  \end{itemize}
\end{itemize}

\subsection*{Relaciones}

\begin{itemize}
  \item Conexión decimal-fraccionario
  \begin{itemize}
    \item \href{https://www.tierradenumeros.com/post/hilo-conexion-decimal-fraccion-tortillas-unpacked/}{Tortilla de patatas}
    \item \href{https://twitter.com/smaccarrone/status/1636353759281307649?s=20}{Número de decimales}
  \end{itemize}
  \item Cuadrados perfectos y raíces cuadradas exactas.
  \item Factores, múltiplos y divisores. Factorización de números primos.
  \begin{itemize}
    \item Descomposición en línea. (Buscar los divisores de números, y dejar que surjan diferentes estrategias)
    \item Actividades de búsqueda de patrones y relaciones en la tabla del cien (Ruiz-López, 2000 y 2003)
    \item Divisibilidad
    \begin{itemize}
      \item \href{https://nrich.maths.org/6651}{Aritmética Modular}
    \end{itemize}
    \item mcm y mcd
    \begin{itemize}
      \item \href{https://nrich.maths.org/mobile?utm_source=secondary-map}{Actividades nrich}
    \end{itemize}
  \end{itemize}
  \item Comparación y orden de fracciones: exacta y aproximada.
\end{itemize}

\subsection*{Proporcionalidad}

\begin{itemize}
  \item Razón y proporción: magnitudes directamente proporcionales.
  \item Porcentajes: Aumento y disminución porcentual.
  \item Situaciones de prpoprcionalidad directa e inversa. Análisis y estrategias.
\end{itemize}

\subsection*{Educación Financiera}

\begin{itemize}
  \item Toma de decisiones basada en calidad-precio y valor-precio.
  \begin{itemize}
    \item Trabajar a través de los porcentajes, para comparación de precios.
  \end{itemize}
\end{itemize}

\section*{(B2) Medida y geometría}

\subsection*{Magnitud}

\subsection*{Medición}

\subsection*{Estimación y funciones}

\section*{(B3) Geometría en el plano y el espacio}

\subsection*{Figuras geométricas 2D}

\subsection*{Localización y sistemas de representación}

\section*{(B4) Figuras geométricas 2D}

\section*{(B5) Álgebra}

\section*{(B6) Estadística}
